\documentclass{article}
\usepackage{amsmath}
\usepackage{hyperref}
\usepackage{geometry}
\geometry{margin=1in}
\begin{document}

\section*{Benchmark Methods: Technical Details}

This section provides a brief technical overview of the polynomial root-finding methods used in the benchmarks: the companion matrix method from the \texttt{Polynomials.jl} package, and the homotopy continuation method from \texttt{HomotopyContinuation.jl}. This information is intended to help the reader understand the underlying algorithms and their typical use cases.

\subsection*{1. Polynomials.jl: roots}
\textbf{Package:} \href{https://github.com/JuliaMath/Polynomials.jl}{Polynomials.jl}\\
\textbf{Function:} \texttt{roots(p::Polynomial)}

The \texttt{roots} function in \texttt{Polynomials.jl} computes all roots (zeros) of a univariate polynomial with floating-point coefficients. The default method is based on the \emph{companion matrix} approach:
\begin{itemize}
  \item The polynomial $p(x) = a_n x^n + \cdots + a_0$ is represented by its companion matrix $C$.
  \item The eigenvalues of $C$ correspond to the roots of $p(x)$.
  \item The eigenvalue computation is performed using LAPACK routines, which are highly optimized for dense matrices.
\end{itemize}
This method is efficient for polynomials of moderate degree (up to a few hundreds), but may suffer from numerical instability for very high degrees or ill-conditioned polynomials. The implementation is located in the file \texttt{src/roots.jl} of the package.

\subsection*{2. HomotopyContinuation.jl: solve}
\textbf{Package:} \href{https://github.com/JuliaHomotopyContinuation/HomotopyContinuation.jl}{HomotopyContinuation.jl}\\
\textbf{Function:} \texttt{solve(F; kwargs...)}

The \texttt{solve} function in \texttt{HomotopyContinuation.jl} computes all isolated solutions of a system of polynomial equations using \emph{numerical homotopy continuation}:
\begin{itemize}
  \item The method constructs a homotopy $H(x, t)$ that continuously deforms a system with known solutions into the target system $F(x)$ as $t$ goes from 0 to 1.
  \item Solution paths are tracked numerically from the start system to the target system using predictor-corrector methods.
  \item The algorithm is robust for both univariate and multivariate systems, and can handle singularities and positive-dimensional solution sets.
\end{itemize}
Homotopy continuation is a general and powerful approach, especially for systems where symbolic methods are infeasible. The implementation is located in the directory \texttt{src/} of the package, with core routines in \texttt{src/solve.jl} and \texttt{src/pathtracker.jl}.

\paragraph{References}

\subsection*{3. PolynomialCovers.jl: cyclic\_lift\_solver}
\textbf{Package:} PolynomialCovers.jl\\
\textbf{Function:} \texttt{cyclic\_lift\_solver(P, y; s=3, tol=1e-9, ...)}

The \texttt{cyclic\_lift\_solver} function computes the roots of $P(x) = y$ by exploiting the additional structure of polynomial covers. The method is based on fast monodromy lifting from a base fiber, under the assumption that $|y| > 3S/2$, where $S$ bounds the critical value disk of the cover:
\begin{itemize}
  \item The algorithm constructs a base fiber using monodromy at infinity, ensuring all roots are found in a numerically stable region.
  \item Each root is tracked independently along a path in the complex plane, using path-lifting techniques.
  \item The roots are returned in their natural cyclic order, reflecting the monodromy structure of the cover.
\end{itemize}
\textbf{Numerical advantages:}
\begin{itemize}
  \item The method is highly stable and efficient when the target value $y$ is outside the critical disk, as all paths are well-separated and avoid singularities.
  \item Parallelization is straightforward, since each root can be lifted independently.
  \item The cyclic ordering of the roots is preserved, which is valuable for applications where the monodromy or sheet structure is relevant (e.g., algebraic curves, covering spaces).
\end{itemize}
The implementation is located in \texttt{src/PolynomialCovers.jl}, and relies on auxiliary routines for fiber construction and path lifting.

\textbf{Reference:} \url{https://github.com/luisjavierhernandez/PolynomialCovers}

\paragraph{References}
\begin{itemize}
  \item Polynomials.jl: \url{https://github.com/JuliaMath/Polynomials.jl}
  \item HomotopyContinuation.jl: \url{https://github.com/JuliaHomotopyContinuation/HomotopyContinuation.jl}
  \item PolynomialCovers.jl: \url{https://github.com/luisjavierhernandez/PolynomialCovers}
  \item Trefethen, L. N., & Bau, D. (1997). \emph{Numerical Linear Algebra}. SIAM. (Companion matrix method)
  \item Sommese, A. J., & Wampler, C. W. (2005). \emph{The Numerical Solution of Systems of Polynomials Arising in Engineering and Science}. World Scientific. (Homotopy continuation)
\end{itemize}



\end{document}
